%! AP Statistics — Unit 5, Lesson 5.4 — Homework
\documentclass[10pt]{article}
\usepackage{apstats-article}
\usepackage{apstats-boxes}
\usepackage{grading}
\gradingassignment{apstats-u05-l04-hw}
\gradingcourse{advanced_statistics}

\begin{document}

\pageheader{Unit 5, Lesson 5.4}{Homework}
\namedateperiod

\vspace{0.15in}

\begin{enumerate}

  \item For each statistic, state whether it is an \emph{unbiased} or \emph{biased}
        estimator of the corresponding population parameter. If biased, state the direction.
        \begin{enumerate}[label=\alph*.]
          \item Sample mean $\bar{x}$ as an estimator of $\mu$.
                \gradedhide{\writelines{1}}\gradedright{\gradedblank[id=hw1a, type=short, points=1, answer=unbiased, check=text]{2.2in}{0.17in}}
          \item Sample proportion $\hat{p}$ as an estimator of $p$.
                \gradedhide{\writelines{1}}\gradedright{\gradedblank[id=hw1b, type=short, points=1, answer=unbiased, check=text]{2.2in}{0.17in}}
          \item Sample minimum as an estimator of population minimum.
                \gradedhide{\writelines{1}}\gradedright{\gradedblank[id=hw1c, type=short, points=1, answer={biased high}, check=text]{2.2in}{0.17in}}
          \item Sample variance $s^2$ (using $n-1$) as an estimator of $\sigma^2$.
                \gradedhide{\writelines{1}}\gradedright{\gradedblank[id=hw1d, type=short, points=1, answer=unbiased, check=text]{2.2in}{0.17in}}
        \end{enumerate}

  \item A population of daily step counts has $\mu = 8{,}500$ steps and $\sigma = 2{,}200$.
        \begin{enumerate}[label=\alph*.]
          \item Find $\sigma_{\bar{x}}$ for $n = 121$.\gradedbelow[r]{0.5in-0.32in}{$\sigma_{\bar{x}} = $~\gradedblank[id=hw2a, type=short, points=1, answer=200, check=numeric, tolerance=0]{0.9in}{0.27in}}
                \vspace{0.5in}
          \item How does $\sigma_{\bar{x}}$ change if $n$ increases from 121 to 484?
                Show calculations.\gradedbelow[r]{0.6in-0.32in}{$\sigma_{\bar{x}} = $~\gradedblank[id=hw2b, type=short, points=2, answer=100, check=numeric, tolerance=0]{0.9in}{0.27in}}
                \vspace{0.6in}
          \item Interpret the change in part (b) in a sentence.
                \gradedhide{\writeline\\\gradedhere{\gradedblank[id=hw2c, type=rubric, points=2, display={Quadrupling $n$ halves the SD of $\bar{x}$, so the estimate of mean daily steps is twice as precise.}]{\linewidth+6pt}{0.2in}}\writelines{1}}
        \end{enumerate}

  \item A simulation repeatedly takes samples of size $n = 50$ from a right-skewed
        population with $\mu = 30$ and $\sigma = 8$ and records $\bar{x}$ each time.
        After 1000 trials, the average of all the $\bar{x}$ values is 30.1.
        \begin{enumerate}[label=\alph*.]
          \item Is this evidence that $\bar{x}$ is biased? Explain why or why not.
                \gradedhide{\writeline\\\gradedhere{\gradedblank[id=hw3a, type=rubric, points=2, display={No: the 0.1 difference from 30 is just sampling variability, not evidence of bias.}]{\linewidth+6pt}{\gradedlinesheight{1}}}\writelines{1}}
          \item What would you expect the average of all $\bar{x}$ values to equal
                in the long run?
                \gradedhide{\writelines{1}}\gradedright{\gradedblank[id=hw3b, type=short, points=1, answer=30, check=numeric, tolerance=0]{1.1in}{0.17in}}
        \end{enumerate}

\end{enumerate}

\begin{extensionbox}
\textbf{Extension (optional).}
Two estimators of the same parameter both have $\sigma = 5$. Estimator A is unbiased
($E(A) = \theta$). Estimator B is biased: $E(B) = \theta + 2$. Compute the mean
squared error (MSE = bias$^2$ + variance) for each. Which has smaller MSE?
\end{extensionbox}

\begin{reflectionbox}
\textbf{Preview:} Lesson 5.5 applies what you know about unbiasedness and the CLT to the
sample proportion $\hat{p}$. You'll derive the center, spread, and shape of the
sampling distribution of $\hat{p}$.
\end{reflectionbox}

\end{document}
